\documentclass{article}
\usepackage{pathfinder-readable}

\title{Risk-aware allocation from sampled losses: a static connection between Q and P}
\pathfinderpapers{2608.29130}{A Systematic Approach to Mechanism Design with Stochastic Dynamic
Stability}{2609.04460}{Distributed risk-averse optimization via CVaR}

\begin{document}
\maketitle

\begin{abstract}
Q designs payments for a strategic resource allocation game, while P studies how a network can optimise conditional
value at risk (CVaR) from sampled losses. The thread asked whether P's local risk estimate could be used in Q's game
while retaining efficient allocation, balanced payments and a participation guarantee. The peers found that Q's
original payment can fail the participation guarantee, but a payment correction based only on other agents' messages
repairs it without changing equilibrium allocations. P's estimate then gives bounds on welfare loss and incentives for
the true risk objective, provided the sampling queries and a planner multiplier exist. The verifier ended the thread at
DRAFT because the corrected static argument survived a general-agent check; convergence of sampled play in Q's game
remains unproved.
\end{abstract}

\section{Introduction}

Q considers several agents who each request an allocation from a shared, limited supply. Each agent values its own
allocation, and a manager uses a quadratic payment rule to make the agents' choices agree with the allocation that
maximises total value. Q also proposes a sampled best-response method and claims a unique equilibrium and convergence
under its stated conditions \cite{Q}.

P considers agents on a changing communication network who all seek the same decision. Their losses are uncertain, and
they minimise the average of their own CVaR objectives, a measure of severe losses in the upper tail. They can sample
losses but cannot evaluate the risk objective or its gradient directly. P constructs empirical CVaR estimates, uses
them in a network update, and proves consensus and optimisation results for that shared-decision problem \cite{P}.

The thread asked whether P's sampled risk values could replace Q's valuations. The objects that agents choose differ: Q
needs separate allocations, often unequal, while P brings copies of one decision into agreement. The peers therefore
used P's local estimate, not its consensus iteration. They examined Q's payment rule before making that connection.

\section{What was done}

The peers first checked Q's claimed strict uniqueness condition. Its implementation condition makes the payment's price
response unchanged when all agents raise their price messages together. Along that common-price direction, the
quadratic form used for strict uniqueness is zero. It therefore cannot be strictly negative as Q requires. This
objection did not decide whether Q's allocation rule works, so the peers checked its displayed payment directly.

Let there be \(N\geq 2\) agents. Agent \(n\) chooses an allocation \(x_n\) from a compact convex set \(\mathcal X_n\)
and a nonnegative price message \(p_n\). The resource capacity is the nonnegative vector \(c\), and \(V_n(x_n)\) is
agent \(n\)'s valuation. Q's payment is \(t_n\), and its positive scaling parameter is \(\alpha\). For any one
resource, write \(z\) for total requested allocation minus capacity and \(P\) for the sum of price messages. The peers
calculated the derivative of an agent's payoff with respect to its price message:
\[
\frac{\partial (V_n-t_n)}{\partial p_n}
=-\frac{\alpha}{N-1}(Np_n-P-z).
\]
If any equilibrium price for that resource is positive, all prices there are equal and \(z=0\). If all are zero,
\(z\leq 0\). Applied to every resource, this gives a common nonnegative price vector \(p\), feasible allocations and
\(p^\top(\sum_n x_n-c)=0\). At that price, each agent's allocation choice maximises \(V_n(x_n)-\alpha p^\top x_n\).
These are the planner's allocation conditions with multiplier \(\lambda=\alpha p\). Under strong concavity, the planner
allocation is unique; a multiplier need not be. Conversely, if a planner multiplier exists, the peers' joint-concavity
check shows that a planner solution with its common price is an equilibrium.

They then tested Q's claimed participation guarantee. Take \(N=2\), a single resource with capacity \(1\), allocation
sets \([0,1]\), \(\alpha=1\), and valuations \(V_1(x)=4x-x^2/2\) and \(V_2(x)=x-x^2/2\). The allocation \((1,0)\) and
common price \(3\) form an equilibrium. Q's payment charges the first agent \(9/2\), more than its value \(7/2\). Its
payoff is \(-1\), below the zero-allocation option. A separate example with \(N=2\) has the unique allocation \((1,1)\)
but every common price from \(0\) to \(1\) as an equilibrium price; it confirms that message uniqueness fails even
where allocation is fixed.

The peers next added a term to Q's payment that uses only other agents' messages. Let \(\bar p_{-n}=\sum_{m\ne n}p_m\)
and \(\bar x_{-n}=\sum_{m\ne n}x_m\). Their added term is
\[
g_n=\frac{\alpha N}{(N-1)^3}\bar p_{-n}^{\top}
\left(\bar x_{-n}-\frac{N-1}{N}c\right).
\]
The revised payment is \(t_n+g_n\). Since \(g_n\) does not depend on agent \(n\)'s messages, it changes no best
response. Finally, the peers used P's expected empirical CVaR, averaged over fresh samples and small random
perturbations, as a fixed valuation surrogate. They compared its equilibria with the planner problem built from true
CVaR.

\section{Results}

The payment calculation matters for any number of agents. Put \(d_n=x_n-c/N\), \(z=\sum_m x_m-c\), and
\(\kappa_N=1+N/(N-1)^2\). At a common price, direct expansion of Q's payment gives
\[
\frac{t_n}{\alpha}=\kappa_Np^\top d_n
-\frac{N}{(N-1)^2}p^\top z,
\qquad
\frac{g_n}{\alpha}=\frac{N}{(N-1)^2}p^\top(z-d_n).
\]
Thus the revised payment is \(\lambda^\top(x_n-c/N)\) at every common-price profile. At equilibrium, the payments sum
to zero by capacity complementarity. If \(0\in\mathcal X_n\) and valuations are normalised so that \(V_n(0)=0\), the
agent's allocation choice gives
\[
V_n(x_n)-\lambda^\top(x_n-c/N)
\geq \lambda^\top c/N\geq 0.
\]
The correction therefore restores budget balance and participation at every equilibrium while keeping Q's efficient
allocation. Emmy's transfer audit in \texttt{../emmy/transfer\_audit.md} contains the full expansion and a three-agent
check.

For the risk connection, let \(J_n(x,\xi)\) be agent \(n\)'s sampled loss at allocation \(x\) under uncertainty
\(\xi\), and let \(C_n(x)\) be its true CVaR at upper-tail level \(\beta_n\). Let \(H_n(x)\) be the expected empirical
CVaR from a fresh batch of \(s_n\) losses, smoothed by a uniform perturbation in a ball of radius \(\delta_n\). Suppose
the losses are convex, \(L_n\)-Lipschitz and bounded in magnitude by \(U_n\) throughout the queried neighbourhood. P's
CVaR and empirical-distribution estimates, together with its smoothing bound, give the pointwise comparison
\[
|H_n(x)-C_n(x)|\leq h_n,
\qquad
h_n=L_n\delta_n+\frac{U_n\sqrt{2\pi}}{\beta_n\sqrt{s_n}}.
\]
This is a bound on a deterministic expected surrogate at each point, not a claim that every random batch has this error.

To meet Q's curvature condition, let \(u_n(x)\) be a strongly concave utility term with strength \(\alpha\). Define the
true valuation from \(u_n-C_n\), and the surrogate valuation from \(u_n-H_n\), subtracting each one's value at \(0\) so
both outside options equal zero. With feasible perturbation queries and a multiplier for the surrogate planner problem,
the corrected game has an equilibrium. Each such equilibrium maximises surrogate welfare and has exactly balanced
payments and nonnegative surrogate payoff. Its price messages may differ across equilibria.

Let \(\widetilde x\) be an equilibrium allocation of that surrogate game, and let \(x^\star\) maximise true welfare
over the same capacity-feasible allocation set. The peers' value comparison gives
\[
\text{true welfare at }x^\star-\text{true welfare at }\widetilde x
\leq 2\sum_n h_n,
\qquad
\|\widetilde x-x^\star\|\leq
2\sqrt{\frac{\sum_n h_n}{\alpha}}.
\]
For agent \(n\), the gain from any unilateral change of message when payoffs use true valuations is at most \(2h_n\),
and its true payoff is at least \(-2h_n\) relative to its normalised outside option. The revised payments still balance
exactly. The welfare and incentive bounds compare two decisions, so the constants introduced by normalising valuations
cancel. Strong concavity supplies the distance bound. Ada's derivation in \texttt{../ada/pair\_note.md} and Emmy's
independent audit support these static conclusions.

\section{Objections and limits}

Q's strict uniqueness condition conflicts with its implementation condition, and its original payment can violate
participation. The direct equilibrium argument and the others-only correction settle allocation efficiency, balance and
participation under the stated conditions. They do not make the price unique. The peers also noted that Q's numerical
sample rule equals one sample at every iteration as written, while its theorem requires geometrically increasing
samples. That simulation does not verify the theorem's sampling regime.

P's consensus theorem concerns copies of one shared decision. It does not show that sampled agents reach an equilibrium
of Q's separate-allocation game. P also assumes a feasible set containing a ball around zero; Q's nonnegative
allocation sets do not meet that assumption literally. The local risk estimate needs a neighbourhood on which all
perturbation queries are possible, and the equilibrium argument needs a planner multiplier. These are conditions here,
not established properties of every allocation problem. The objective is the sum of individual CVaRs, not the CVaR of
total loss.

A narrow search by the peers recorded adjacent work on CVaR in a two-stage market \cite{dahlinJain} and learning in
stochastic aggregative games \cite{duMeng}. Those reports were not used in the proofs and do not establish novelty. The
present material remains thin on implementable learning dynamics and wider prior work.

\section{Why the thread stopped}

The verifier first objected that the transfer coefficient was wrong for general \(N\). The peers re-expanded all terms
of Q's displayed payment and checked a three-agent equilibrium. Their result included a term the first objection had
omitted, and the final verifier accepted the corrected coefficient. It marked the thread DRAFT because the static
connection was supported, while the note kept the query and multiplier conditions explicit and made no claim that P's
algorithm converges for Q's game. The smallest step to restart the research line is to prove or refute convergence of a
feasible sampled learning rule to the corrected game's equilibrium set, starting with two agents and one resource.

\bibliographystyle{plainurl}
\bibliography{references}
\end{document}
